\section{A Chow-form determinant and the arithmetic slope}
\label{sec:chow-determinant}
This section gives the determinant estimate used in the component
Hilbert--Samuel inequality. Let \(F\) be the coefficient number field.
We first state exactly the geometric and metric data used in the
calculation. Let \(V_j\), \(0\le j\le T\), be the graded section spaces
with full integral lattices and Hermitian norms at infinity. Put
\(D\ge1,\ T\ge D+1,\ d=\dim_F V_T\). Multiplication by a coordinate
linear form increases the index \(j\) by one. The system has
bases over \(F\) at indices \(0,D,T\) whose vectors lie in the respective lattices and have archimedean norm bounds
\(\Lambda_0,\Lambda_D,\Lambda_T\), all at least one. The complex
multiplication bound for a unit coordinate vector is \(\sqrt\beta\)
with \(\beta\ge1\).

\paragraph{Filtration at an intersection of hyperplanes.}
The projective image admits an irreducible Chow form \(f(s,l)\),
homogeneous of degree \(\delta\ge1\) in each of \(D+1\) hyperplane
blocks. Here \(s=(s_1,\ldots,s_D)\) and \(l\) is the last block.
There is a nonzero polynomial \(Q(s)\) such that at every \(F\)-point
with \(Q(s)\ne0\) the section splits into \(\delta\) distinct reduced
points \(P_1,\ldots,P_\delta\) over a finite extension. Thus
\[
 f(s,l)=c_s\prod_{\nu=1}^\delta l(P_\nu),\qquad c_s\ne0.
\]
\paragraph{The residue map and the subspace it filters.}
Write \(Y\) for the smooth projective \(D\)-fold over the splitting
field, \(L\) for the generated hyperplane bundle and \(A,L_2\) for
the two auxiliary nef and big bundles. Use additive notation for
tensor products. The data include a section
\(\epsilon\in H^0(Y,A-L_2-DL)\) nonvanishing on the prescribed
open \(U\). If \(Z\) is the common zero scheme of \(s_1,\ldots,s_D\),
the good-section conditions are that \(Z\) be finite, \'etale and
contained in \(U\). In this paragraph its \(\delta\) points are
rational after the indicated finite extension. Define
\[
 G_j=\operatorname{im}\{H^0(Y,K_Y+A+jL)\to H^0(U,K_Y+A+jL)\}.
\]
The coefficient space satisfies \(G_j\subset V_j\) for \(j\le T\)
and \(G_0=V_0\). The inclusion can matter for a positive \(j\).

On \(Z\) the ordered sections define the nowhere-zero Jacobian
\(J_s=ds_1\wedge\cdots\wedge ds_D\) in
\((K_Y+DL)|_Z\). Indeed the conormal map is an isomorphism at
each transverse point. In local frames the degree-\(j\) residue
of \(m\in V_j\) is \(m|_Z/(\epsilon J_s)\), with values in
\((L_2+jL)|_Z\). The choice of frames identifies its values with
one scalar at each of the \(\delta\) points.

We show that the degree-zero residue map is onto. Any prescribed
\(v\in H^0(Z,L_2|_Z)\) gives \(J_s v\) in
\(H^0(Z,(K_Y+L_2+DL)|_Z)\). This latter section lifts to
\(H^0(Y,K_Y+L_2+DL)\). To see the lifting, use the Koszul
resolution of \(\mathcal O_Z\) by the \(D\) equations, tensored by
\(K_Y+L_2+DL\). The term of index \(i\) is the direct sum of
\(\binom Di\) copies of \(K_Y+L_2+(D-i)L\), \(0\le i\le D\).
Each has zero higher cohomology by Kawamata--Viehweg vanishing:
\(L_2\) is nef and big, \(L\) is nef, and \(D-i\ge0\). The complex
is exact off degree zero. Indeed a generator is a unit locally away
from \(Z\), and at points of \(Z\) the equations are a regular system
of parameters. The cohomology of its global sections therefore computes
that of \((K_Y+L_2+DL)|_Z\). This gives the asserted lift \(b\).
Then \(\epsilon b\) restricts to an element of \(G_0=V_0\) whose
residue is \(v\). Equivalently one may first change the equations by
an invertible constant matrix to a sequence regular on successive
quotients, as in the source's restriction argument.

Let \(I\) be the ideal of \(Z\). A first-order vanishing section
\(g\in G_{l+1}\) with \(D\le l+1\) can be written
\[
                        g=\sum_{i=1}^D s_i g_i,\qquad g_i\in G_l .
\]
Its global extension restricts to zero on \(Z\), since \(Z\subset U\).
Apply the same Koszul complex tensored by \(K_Y+A+(l+1)L\).
Its terms have \(L\)-exponents \(l+1-i\ge0\), and the higher
cohomology again vanishes. Exactness of global sections at its
index-zero term identifies the restriction kernel with the image
of the \(D\) indicated multiplications. This proves the division.

We need the version for higher orders with an explicit lower bound
on the remaining degree. For \(j\ge D\), \(a\ge0\), every
\(g\in G_{j+a}\) vanishing to order \(a\) along \(Z\) admits
\[
                       g=\sum_{|\alpha|=a}s^\alpha g_\alpha,
                       \qquad g_\alpha\in G_j .
\]
For \(a=0\) this is immediate. For order \(a+1\), apply the
order-\(a\) result with coefficient degree \(j+1\).
Locally along \(Z\) the associated graded algebra of \(I\) is the
polynomial algebra on the conormal basis given by the \(s_i\);
this is the regular-sequence description
\(\operatorname{gr}_I\mathcal O_Y=\operatorname{Sym}_{\mathcal O_Z}
(I/I^2)\). Thus the degree-\(a\) symbol of the resulting expression
has a separate coefficient \(g_\alpha|_Z\) for each monomial,
in local frames. Since the section vanishes to order \(a+1\),
every such coefficient is zero. Apply the first-order division
above to each \(g_\alpha\in G_{j+1}\). Regrouping the terms with
the same multi-index \(\alpha+\mathbf e_i\) gives the assertion.

Write \(r=T-D\) and set
\[
 P_s=\sum_{|\alpha|=r}s^\alpha V_D,\qquad W_s=G_T+P_s\subset V_T,
 \qquad U_a=\{h\in W_s:h\text{ vanishes along }I^a\}.
\]
Here vanishing means membership in \(I^a\) in local bundle
coordinates along \(Z\). Stability under multiplication gives the
inclusion \(P_s\subset V_T\). Now \(P_s\subset U_r\) by construction.
Conversely, for \(h=g+p\in U_r\) with \(g\in G_T\), \(p\in P_s\),
the difference \(g=h-p\) vanishes to order \(r\). Higher-order
division with \(j=D\) expresses \(g\) in \(P_s\), proving \(U_r=P_s\).

For \(a<r\), the kernel of the symbol on \(U_a\) is \(U_{a+1}\)
by the same associated graded description. Its target after choosing
frames is \(\bigoplus_{|\alpha|=a} F'^{\,\delta}\) over the splitting
field \(F'\). To show surjectivity, choose \(\ell\) nonvanishing at
all the points, possible over the infinite field \(F'\).
For each \(\alpha\) evaluate the symbol of
\(s^\alpha\ell^{T-a}m\) with \(m\in V_0=G_0\). It vanishes in
every other monomial coordinate, and its \(\alpha\)-coordinate is the
residue vector of \(m\) multiplied at point \(P_\nu\) by
\(\ell(P_\nu)^{T-a}\). These nonzero factors and surjectivity of the
degree-zero residue prove that this block spans its entire target.
All these test sections lie in \(G_T\subset W_s\).
Writing \(y_s=\dim P_s\), summing graded dimensions now gives
\[
            \dim W_s=y_s+\delta\binom{T-1}{D},\qquad
            \dim(V_T/W_s)=d-y_s-\delta\binom{T-1}{D}.
\]
This identifies the exact dimension of the subspace requiring
additional top-basis columns in the determinant. The argument uses
ordinary adjoint sections on the compactification only in the
subspaces \(G_j\), and uses the actual \(V_D\) in \(P_s\).

Let \(R=\binom{T-1}{D}\). For \(0\le a<T-D\) the number of
monomials in \(D\) variables of degree \(a\) is
\(c_a=\binom{a+D-1}{D-1}\). Stars and bars and Pascal summation give
\(\sum c_a=R\). The total exponent of the last form on one residue
point in the determinant below is \(b=\sum_a c_a(T-a)\). Indeed
\(a c_a=D\binom{a+D-1}{D}\), so, taking the sum with \(a<T-D\),
\[
 b=TR-D\binom{T-1}{D+1},\qquad
 (D+1)b=R\bigl(T+D(D+1)\bigr).
\]
These are exact finite identities, including \(T=D+1\).

Choose a rational good \(s_0\) for which the deepest-step dimension
\(y\) is maximal. Such a point exists because \(F\) is infinite, \(Q\)
is nonzero, and the possible dimensions are integers between zero and
\(d\). After splitting the section at \(s_0\), select \(\delta\)
vectors of the short basis of \(V_0\) with independent evaluations.
Choose \(l\) avoiding all the points. For each of the above monomials
\(s^\alpha\) take the \(\delta\) columns \(s^\alpha l^{T-a}m_i\),
where \(m_i\) are the selected degree-zero sections. Their symbols are
independent: at each monomial the invertible evaluation matrix is
multiplied by a diagonal matrix with nonzero entries \(l(P_\nu)^{T-a}\).
The different symbols lie in independent summands or successive
filtration grades. Include \(y\) columns of the form \(s^\alpha m\),
\(|\alpha|=T-D,\ m\) a short degree-\(D\) basis vector, spanning the
deepest step. Complete with \(d-R\delta-y\) vectors from the short top
basis. These choices can be made by selecting from fixed rational
generating sets over the splitting field. After replacing \(s_0,l\)
by variables, their determinant \(P(s,l)\) is therefore a nonzero
polynomial over \(F\).

\paragraph{Divisibility before taking heights.}
At any other good \(s\), the dimension of the deepest step is at most
\(y\). If it is smaller, the \(y\) degree-\(D\) columns are dependent
and the determinant vanishes identically in \(l\). If the dimension
equals \(y\), either those columns are dependent, or they give a basis
of the deepest step. In the latter case take symbols successively
modulo the filtration. Only the \(R\delta\) columns from \(V_0\)
depend on the last block on the successive symbol spaces. The
determinant of their block for monomial \(\alpha\) is a scalar
independent of \(l\) times \(\prod_\nu l(P_\nu)^{T-a}\).
The remaining minors are independent of \(l\). This block determinant
identity holds for all \(l\) by polynomial identity. It includes the
possibility of zero scalar if a specialized residue matrix or final minor
loses rank. Consequently for every good \(s\),
\(P(s,l)=\gamma_s f(s,l)^b\), with \(\gamma_s\) in a splitting field.

To descend the divisibility, choose \(l_0\in F^N\) such that
\(g(s)=f(s,l_0)\) is a nonzero polynomial. Evaluation at \(l_0\)
and at \(l\) gives \(P(s,l)g(s)^b=f(s,l)^bP(s,l_0)\) on \(Q(s)\ne0\).
Multiplying the difference by \(Q\) gives a polynomial vanishing on
all rational tuples. The field is infinite, so the polynomial is zero;
the polynomial ring is a domain, so \(Q\) cancels. The polynomial \(f\)
has positive degree in \(l\) whereas \(g\) is independent of \(l\).
Irreducibility of \(f\) therefore gives \(\gcd(f,g)=1\).
In the polynomial unique factorization domain we obtain
\(P=f^bG\) for a nonzero multihomogeneous \(G\).

\paragraph{The lattice estimate.}
Let \(e\) be the sum of the multidegrees of \(G\). The sums of column
degrees on the two sides of this factorization give
\begin{equation}\label{eq:chow-column-degrees}
 b(D+1)\delta+e=R\delta T+y(T-D),\qquad R\delta+y\le d .
\end{equation}
Write \(H_f,H_G\) for the source's heights of the two forms in the
chosen coordinate frame and \(A_T\) for the arithmetic degree of
the Hermitian top lattice. The normalized multiform height is additive
on products. Its lower bound is minus \(c\) times the sum of block
degrees, with \(c=\log(n\Lambda)+(\log n+1)/2\), where \(n\ge1\)
counts the spanning coordinates and \(\Lambda\ge1\) bounds their
lattice data. Hence \(H_f\ge-c(D+1)\delta\), \(H_G\ge-ce\).
At each finite place the substituted polynomial vectors have integral
lattice coefficients. At each complex place the three types of columns
have norm bounded respectively by
\(\beta^{T/2}\Lambda_0,\ \beta^{(T-D)/2}\Lambda_D,\ \Lambda_T\)
for unit input vectors. The determinant inequality for polynomial
vectors in a Hermitian lattice gives
\begin{align}
 A_T\ \ge\ bH_f+H_G
   &-R\delta\log(\beta^{T/2}\Lambda_0)
       -y\log(\beta^{(T-D)/2}\Lambda_D) \notag\\[-2pt]
   &-(d-R\delta-y)\log\Lambda_T .\label{eq:chow-slope-raw}
\end{align}
This is a statement about the degree of the actual lattice, so possible
finite index of the columns in it is already included. The local
determinant estimate at infinity uses the exterior-product norm
inequality; summing the logarithms with the finite ideal contributions
uses the product formula. The stated multiform height bounds and this
lattice-height comparison are the metric inputs here.

Put \(h_f=H_f/((D+1)\delta)\) and \(z=1-R\delta/d\).
Equation \eqref{eq:chow-column-degrees} shows that \(d>0\),
\(0\le z\le1\) and \(e=y(T-D)-R\delta D(D+1)\le T(d-R\delta)\).
Furthermore
\(bH_f=R\delta (T+D(D+1))h_f
\ge R\delta(Th_f-D(D+1)c)\), because \(h_f\ge-c\).
All the logarithmic losses are nonnegative. In
\eqref{eq:chow-slope-raw} we may bound the loss on the final
\(d-R\delta\) columns by charging each
\(T\log\beta/2+\log\Lambda_D+\log\Lambda_T\). Division by \(d\) gives
\begin{align*}
 A_T/d \ \ge\
 &(1-z)\bigl(Th_f-D(D+1)c-\log\Lambda_0-(T/2)\log\beta\bigr)\\
 &-z\bigl(Tc+(T/2)\log\beta+\log\Lambda_D+\log\Lambda_T\bigr).
\end{align*}
This is the intermediate slope inequality in the source.

\paragraph{Comparison with the component height.}
The frame comparison gives \(h_f\ge h-(\log n+1)/2\), where \(h\)
is the geometric height used in Section~\ref{sec:hs-component}. Set
\[
 L=\log\Lambda+26\log\!\bigl(nD\delta(2+|h|+\log\Lambda)\bigr).
\]
Since \(nD\delta(2+|h|+\log\Lambda)\ge2n\), we have \(c\le L\).
In the first bracket of the slope inequality we may replace
\(Th_f\) by \(Th-T(\log n+1)/2\).
The three bounds
\[
 D(D+1)c\le26D^2(L+\log T),\qquad
 (T/2)\log\beta\le(T+D+1)\log\beta,
\]
\[
 \frac{T(\log n+1)}2
 \le (T+D^2)\left(\frac{(D+1)(\log n+1)}2+3\right)
\]
allow the additional subtraction \(3(D+1)|h|\) used by the component
formula. In the second bracket replace \(c\) by \(L+\log2\).
Both replacements weaken a lower bound since \(z,1-z\ge0\).
With \(J=\log\Lambda_0\) these are exactly the two terms of
\srcid{HSClaimS} at constant \(26\). The inputs on the geometric side are Bertini for the simple hyperplane
sections, adjoint vanishing on the smooth projective model and the
short-basis assertions above. Identifying the
birational image degree, its Chow height and the Hermitian lattice
degree transfers the inequality to the component of the window.

\paragraph{Source correspondences.}
The restriction and filtration are
\srcid{HilbertSamuelDegreeZeroResiduesRangeEqTop},
\srcid{FilMonoOfKerRes}, \srcid{FilSpaceVEqSpaceP} and
\srcid{HilbertSamuelVanishingFiltrationDeepestStepSpan}.
The selection of columns and the specialization identity are
\srcid{ExistsShortColumnsDetNeZeroHilbertSamuel} and
\srcid{DetRationalSectionsEqPowProductLinearForms}. The descent is
\srcid{PowDvdOfSpecialisations}; the weighted count and its numerical
use are \srcid{MonomialCountsDegreeAndWeightedSum} and
\srcid{ClaimSOfNumbers}. See also \srcid{SlopeLowerBoundTopLatticeOffHypersurface},
\srcid{ExplicitArithmeticHilbertSamuelBound}, and
\srcid{HilbertSamuelComponentShortBasisBound}.
